\section{共鳴の規約と入力の範囲}
Rainsの既知のeffective-root理論，境界triangleと反標準環の構成を入力とする．
zigzagの形式性はLiu--Wangの既知定理であり，本研究の新定理として扱わない．
全$q$の導来比較は主論文\cite{Hearts}の\texttt{dp-hearts.tex}に証明がある．
主論文ではさらに，全ADE境界外heartのpreprojective模型と，
元の$A$の全有限次元表現を支配する商を扱う．本稿は符号別Ext・混合拡大・形式変形と，三次曲面のcornerの完全計算を主題とする．

基礎体を$k=\mathbb C$とする．Rainsの意味のsmooth elliptic
noncommutative del Pezzo $(\mathcal A_{\mathrm R},E)$に
構造line bundle $e=\mathcal O_X$，Serre関手$\mathbb S$，
$\theta=\mathbb S[-2]$，$\tau=\theta^{-1}$，
境界制限$Li^*:\Db(\mathcal A_{\mathrm R})\to\Perf(E)$
と右随伴$i_*$があり，
\begin{equation}\label{eq:boundarytriangle}
 \theta M\longrightarrow M\longrightarrow i_*Li^*M
 \longrightarrow\theta M[1]
\end{equation}
は自然なtriangleである\cite[\S\S4--5]{RainsFiltered}．
$e$はline bundleで$\theta,\tau$はRains heartにexactに作用する．
さらに$R=\bigoplus_{n\ge0}\Hom(e,\tau^ne)$は
次数1の中心元$g$を持ち，$R/(g)$は楕円曲線上のtwisted
homogeneous coordinate ringとなる\cite[Proposition 5.1]{RainsFiltered}．
以下，$\qgr R$の対象と有限生成次数付き加群を扱う場合には，
当該$R$の右noetherian性と必要な$\chi$条件を仮定または確認した範囲に限る．

古典的avatarを滑らかなgenuine del Pezzo曲面$S$，
その滑らかな反標準曲線を$E$とし，
$K_0(\mathcal A_{\mathrm R})\cong K_0(S)$を固定する．
根$\alpha\in K_S^\perp$は$\alpha^2=-2$を満たす．
$M$のrankは零，$c_1(M)=\alpha$，
$m=\chi(e,M)$とする．
gluingパラメータ$q\in\Pic^0(E)$の向きを
\begin{equation}\label{eq:resonance}
 Li^*M=0\quad\Longrightarrow\quad
 \mathcal O_E(\alpha)\otimes q^{m}\simeq\mathcal O_E
\end{equation}
と規格化する．これは$K$群上のrestriction determinant
の式である．Rainsの$q$を逆元とする規約では指数の符号も反転する．
式\eqref{eq:resonance}だけではsheafの存在を保証しない．
以下では，Rainsの既約根sheaf存在定理
\cite[\S9, \S11, pp.~139--141]{RainsBirational}と
\cite[\S4, pp.~33--35]{RainsFiltered}が適用され，
$Li^*M=0$である\emph{既約な純一次元sheaf}が実在する共鳴を扱う．
特に，固定したblowdownのsimple rootの共鳴を使えばよい．
非torsion仮定を使う単純性・有限長性は\S\ref{sec:mixed-a2-category}で別に証明する．

\begin{proposition}[共鳴による反標準環のHilbert関数の跳びはない]
\label{prop:hilbert-no-jump}
前節のRainsの環の仮定の下，$d=(-K_S)^2>0$とすると
\[
 \dim_k R_n=1+\frac d2 n(n+1)\qquad(n\ge0)
\]
であり，$q$の共鳴の有無に依存しない．
特に非零指数の共鳴で生じる有限次元アフィン表現は，
$R$のHilbert関数の跳びと同一視してはならない．
\end{proposition}
\begin{proof}
$g$は正則中心元であり
$R/(g)=B(E,L,\sigma)$で$\deg L=d$だから
\[
0\longrightarrow R_{n-1}\xrightarrow{\cdot g}R_n
\longrightarrow B(E,L,\sigma)_n\longrightarrow0
\]
が成立する．$n\ge1$では楕円曲線上の正次数線束の
切断数は$nd$である．$R_0=k$から帰納法で主張が従う．
このgraded ringの恒等式はRains
\cite[Proposition 5.1]{RainsFiltered}の仮定内での結果である．
\end{proof}

\section{根sheafのEuler指数とcohomological degree}
\label{sec:sign-trichotomy}
Rainsの$\mathrm{ad}$は，一般に同じ非可換曲面の自己双対ではなく，
反対側の双対曲面$X^{\mathrm{ad}}$への反変導来同値である．
構造line bundle $e$はその双対$e^{\mathrm{ad}}=\mathrm{ad}(e)$へ移る．
既知の双対性
\[
 \mathbf R\mathrm{ad}_X:
 D^b(\mathcal A_{\mathrm R})^{\mathrm{op}}
 \xrightarrow{\sim}D^b(\mathcal A_{\mathrm R}^{\mathrm{ad}})
\]
とその二重双対を使う\cite[\S5, Corollary 5.9]{RainsBirational}．

\begin{lemma}[純一次元の双対と指数反転]
\label{lem:root-dual}
$M$を境界に制限して零となる既約な純一次元sheafとする．
$N=R^1\mathrm{ad}_X(M)$は双対曲面上の既約な純一次元sheafで，
\[
 \mathbf R\mathrm{ad}_X M\simeq N[-1],\quad
 R^1\mathrm{ad}_{X^{\mathrm{ad}}}N\simeq M,\quad
 c_1(N)=c_1(M),\quad \chi(e^{\mathrm{ad}},N)=-\chi(e,M).
\]
$M$上で自然な$\theta M\to M$が同型なら，
$N$上でも自然な$\theta^{\mathrm{ad}}N\to N$が同型である．
とくに，根sheafは双対でも根sheafになる．
\end{lemma}
\begin{proof}
Rains\cite[Corollaries 7.26, 7.30, 7.31]{RainsBirational}
によれば純一次元sheafでは$\mathrm{ad}$の導来cohomologyが
次数1だけにあり，その二重双対は元のsheafを復元する．
これにより$N$も純一次元である．
既約性はgenericな一次元support上の単純加群性に関する性質であり，
genericでの双対によって保たれる．
自然な境界triangleの双対は，双対側の境界triangleに一致するため，
自然変換が$M$上で可逆なら$N$上でも可逆である．
K群では$[\mathbf R\mathrm{ad}M]=-[N]$である．
Rains\cite[Proposition 7.16]{RainsBirational}の公式
$c_1(\mathbf R\mathrm{ad}M)=-c_1(M)$，
$\chi(\mathbf R\mathrm{ad}M)=\chi(M)$を使えば
$c_1(N)=c_1(M)$，$\chi(N)=-\chi(M)$となる．
ここで$\chi(N)=\chi(e^{\mathrm{ad}},N)$である．
\end{proof}

\begin{lemma}[零以下のEuler指数ではglobal sectionがない]
\label{lem:no-section}
$M$を根$\alpha^2=-2$を第一Chern類とする既約な純一次元sheafとする．
$m=\chi(e,M)\le0$なら$\Hom(e,M)=0$．
\end{lemma}
\begin{proof}
非零の$f:e\to M$があるとし，その像を$I$とする．
$M$の既約性は「任意の非零部分sheafが同じ第一Chern類を持つ」
ことであって，$f$が全射であることを意味しない
\cite[\S11, pp.~138--140]{RainsBirational}．
$I\ne0$だから$c_1(I)=\alpha$であり，商$T=M/I$は
零または零次元である．
$I$は$e$の商であるため，
Rains\cite[Corollary 10.16]{RainsBirational}の不等式
\[
 \chi(e,I)\ge-\tfrac12\alpha(\alpha+K_X)=1
\]
が適用できる．
一方，零次元sheafのEuler標数は非負
\cite[Proposition 7.6]{RainsBirational}なので
$\chi(e,M)=\chi(e,I)+\chi(e,T)\ge1$となり，
$m\le0$に矛盾する．
\end{proof}

\begin{theorem}[resonance sheafの符号別Extの完全判定]
\label{thm:sign}
$M$は$\alpha^2=-2$，$(-K_X)\alpha=0$を満たす根の
既約な純一次元sheafで，$Li^*M=0$とする．
$m=\chi(e,M)\in\mathbb Z$について全ての$n\in\mathbb Z$で
\[
\boxed{
\begin{aligned}
\dim\Hom(e,\tau^nM)&=\max(m,0),\\
\dim\Ext^1(e,\tau^nM)&=\max(-m,0),\\
\Ext^2(e,\tau^nM)&=0.
\end{aligned}}
\]
したがって自然な導来切断を$M$に適用すると，
\begin{equation}\label{eq:trichotomy}
 \Phi(M)=
 \begin{cases}
   \pi N_m,&m>0,\\
   0,&m=0,\\
   (\pi N_{-m})[-1],&m<0,
 \end{cases}
\end{equation}
ただし$m\ne0$のとき$\pi N_{|m|}\in\qgr R$は
境界$g=0$の外にある非零対象であり，
$R/(g-1)$の$|m|$次元右加群に対応する．
$m=0$のとき，$M$は非零であるが，全Serre orbitに
導来的に右直交し，指定した反標準導来切断は$M$を零に送る．
\end{theorem}
\begin{proof}
$Li^*M=0$なら境界triangleから自然な
$\theta M\to M$が同型であり，$\tau^nM\simeq M$．
line bundle $e$と一次元sheafに対する
$\Ext^2(e,M)=0$は
Rains\cite[Corollary 7.23]{RainsBirational}による．

$m\le0$では補題\ref{lem:no-section}により
$\Hom(e,M)=0$，Euler標数より$\dim\Ext^1(e,M)=-m$．
特に$m=0$では全$\Ext$が零となる．

$m>0$とする．双対$N=R^1\mathrm{ad}M$は
補題\ref{lem:root-dual}により双対曲面上の
負指数$-m<0$の既約根sheafで，境界にも制限して零である．
従って$\Hom_{X^{\mathrm{ad}}}(e^{\mathrm{ad}},N)=0$，
$\dim\Ext^1_{X^{\mathrm{ad}}}(e^{\mathrm{ad}},N)=m$．
Rainsの反変導来双対とSerre双対性を合成すると
\begin{equation}\label{eq:dual-ext}
\begin{aligned}
\Hom_X(e,M)
 &\cong\Ext^1_{X^{\mathrm{ad}}}(e^{\mathrm{ad}},N)^*,\\
\Ext^1_X(e,M)
 &\cong\Hom_{X^{\mathrm{ad}}}(e^{\mathrm{ad}},N)^*.
\end{aligned}
\end{equation}
実際$\mathbf R\mathrm{ad}_XM=N[-1]$なので
\[
 \Ext^{j}_{X^{\mathrm{ad}}}(e^{\mathrm{ad}},N)
 \cong\Ext_X^{j+1}(M,e)
 \cong\Ext_X^{1-j}(e,\theta M)^*
\quad (j=0,1).
\]
これに$\theta M\simeq M$を代入すれば
$\dim\Hom_X(e,M)=m$，$\Ext_X^1(e,M)=0$を得る．

最後に$g$は$\theta\to\mathrm{id}$から生じる中心正則元で，
$M$上で同型として作用するので，
各次数の$\Hom$または$\Ext^1$の$g$による移行は同型．
非零の次数付き$R$加群$N_{|m|}$は
$k[g]$上有限生成で，各次数の次元が$|m|$，
局所化の次数0部分は$|m|$次元の$R/(g-1)$右加群になる．
非零cohomologyが一つの次数に集中しているので
\eqref{eq:trichotomy}が従う．
\end{proof}

\begin{proposition}[root sheafは2-spherical]
\label{prop:spherical}
定理\ref{thm:sign}の$M$は任意の$m$について
\[
\Ext^j(M,M)=
\begin{cases}k&j=0,2,\\0&j\ne0,2.\end{cases}
\]
特に$\theta M\simeq M$の下で$2$-spherical objectである．
\end{proposition}
\begin{proof}
$M$の非零自己準同型の像は既約性により
第一Chern類$\alpha$を持つ．
核が非零なら零次元であり，$M$の純一次元性に反する．
よって自己準同型は単射であり，その余核は零次元．
自己射なので余核の数値類は零，Rains
\cite[Proposition 7.6]{RainsBirational}から余核も零．
従って$\End(M)$は有限次元の斜体で，$k=\mathbb C$だから
$\End(M)=k$．Serre双対により
$\Ext^2(M,M)=\End(M)^*=k$．
Rains\cite[Corollary 7.13]{RainsBirational}のEuler形式で
$\chi(M,M)=-\alpha^2=2$となるので$\Ext^1(M,M)=0$．
\end{proof}

\begin{corollary}[pointedなheart非一致と導来切断の限界]
\label{cor:pointed-sign}
$m<0$の既約根sheafがあると，Rains heartの$(e,\tau)$は
AZ-ampleでない．標準pair$(\pi R,(1))$がAZ-ampleなら
$(\mathcal A_{\mathrm R},e,\tau)$と
$(\qgr R,\pi R,(1))$はpointed abel圏として同値でない．
さらに指定した自然な$\Phi$は$m<0$で$t$-exactでなく，
$m=0$で保守的でない．
$m>0$の個々の根sheafの像はheartの次数0にあるが，
そのことだけでは$\Phi$の全体の$t$-exactnessは従わない．
\end{corollary}
\begin{proof}
$m<0$では非零$M$に対して全ての
$\Hom(e,\tau^nM)$が零なので，AZの生成公理が失敗する．
pointedなabel同値はample pairを保存する．
後二者は\eqref{eq:trichotomy}から従う．
\end{proof}
\begin{remark}[先行する証明の訂正と範囲]
「既約root sheafへの非零な構造sheafからの射は全射」
という推論は正しくない．必要なのは
像が同じ第一Chern類を持ち，商が零次元であることだけである．
補題\ref{lem:no-section}がその正しい証明を与える．
$m=0$の対象が実在することは共鳴の数値条件からは従わない．
genuine del Pezzo avatarでは後述の命題により$m=0$の根共鳴はない．
また指定した$\Phi$の非忠実性は，
別の抽象的導来同値の不存在を意味しない．
\end{remark}

\section{境界外のsheafと拡大によるcohomologyの相殺}
\label{sec:offboundary-extensions}
本節は既約根sheafに限らない．既知入力は
Rainsの境界triangle \eqref{eq:boundarytriangle}，
$1$次元以下のsheafへの$\Ext^2$消滅
\cite[Corollary 7.23]{RainsBirational}，および
中心正則元$g$を持つ反標準環の仮定である．
$\mathcal Z$をRains heartの中で
\[
 \mathcal Z=\{M\in\mathcal A_{\mathrm R}:
    \dim M\le1,\ Li^*M\simeq0\}
\]
と置く．$R$は右noetherianであり，必要な$\chi$条件により
$\Phi=\pi\mathbf R\Gamma_*$を該当対象に定義できると仮定する．
$A=R/(g-1)\cong (R[g^{-1}])_0$と書く．
ここで$\pi$は有限次元次数付き加群によるSerre商である．

\begin{theorem}[境界外のsheafに対する二項振幅と長完全列]
\label{thm:offboundary-amplitude}
$\mathcal Z$は拡大に閉じたexact categoryである．
任意の$M\in\mathcal Z$について次が成立する：
\begin{enumerate}
\item 自然な$\theta M\to M$は同型である．従って
$\tau^nM\simeq M$で，全$n\in\mathbb Z$について
$\Ext^j(e,\tau^n M)=0$ $(j\notin\{0,1\})$．
\item
\[
 N^j(M)=\bigoplus_{n\ge0}\Ext^j(e,\tau^nM)
 \qquad(j=0,1)
\]
は次数付き右$R$加群で，$g$による乗法は全次数で同型である．
特に$k[g]$加群として有限生成freeで，そのrankは
$h^j(M)=\dim_k\Ext^j(e,M)$．
\[
 V^j(M)=(N^j(M)[g^{-1}])_0\in\operatorname{mod}_{fd}A,
 \qquad \dim_kV^j(M)=h^j(M)
\]
であり，$H^j(\Phi M)\simeq\pi N^j(M)$．
従って自然な導来切断のcohomologyは次数$0,1$だけで，
どちらも反標準境界$g=0$の外にある．
\item 任意の短完全列$0\to M_1\to M_2\to M_3\to0$
（$M_i\in\mathcal Z$）から，有限次元右$A$加群の
自然な完全列
\begin{equation}\label{eq:offboundary-six}
0\to V^0(M_1)\to V^0(M_2)\to V^0(M_3)
\xrightarrow{\delta} V^1(M_1)
\to V^1(M_2)\to V^1(M_3)\to0
\end{equation}
を得る．この$\delta$はextension classのconnecting morphismで，
$A$-線型である．
\end{enumerate}
\end{theorem}
\begin{proof}
Rainsのdimension filtrationは短完全列に関し安定である．
また$Li^*$はexact triangulated functorなので
両端の制限が零なら中間項の制限も零となる．
従って$\mathcal Z$は拡大に閉じる．

境界triangle \eqref{eq:boundarytriangle}と$Li^*M=0$から
自然な$\theta M\to M$が同型となる．
$\theta$，$\tau$はheart上でexactなので，全twistも
$1$次元以下のsheafとなる．
負のExtは二つのheart対象間で消え，Rainsの
\cite[Corollary 7.23]{RainsBirational}で$\Ext^2$が消える．
この曲面のcohomological dimensionは高々$2$だから，
(1)が従う．

境界triangleの$\theta\to\mathrm{id}$から得られる
次数$1$の中心元$g$は$M$上で同型であり，
$\Ext^j(e,\tau^nM)\to\Ext^j(e,\tau^{n+1}M)$に
同型を誘導する．有限次元性により
$N^j(M)\simeq k[g]\otimes_k\Ext^j(e,M)$
が$k[g]$加群として成り立つ（$R$作用はこの分解で
単純になるとは主張しない）．
ゆえに$\pi N^j$は非零なら境界外の対象で，
dehomogenizationは$A$の有限次元加群$V^j$を与える．
$\Phi M$のcohomologyの同定は$\pi$のexactnessと
$\mathbf R\Gamma_*$の定義による．

短完全列の$\tau^n$-twistに$\RHom(e,-)$を適用すると
cohomologyは次数$0,1$だけなので六項完全列になる．
これを全$n\ge0$で直和し，中心元$g$で局所化した後
次数$0$を取る．局所化と次数$0$の抽出はこの
graded-moduleのクラスでexactであり，
自然性によりconnecting morphismは$A$-線型となる．
\end{proof}

\begin{corollary}[拡大に閉じた正・負の部分圏]
\label{cor:offboundary-sign-exact}
\[
 \mathcal Z_+=\{M\in\mathcal Z:\Ext^1(e,M)=0\},\qquad
 \mathcal Z_-=\{M\in\mathcal Z:\Hom(e,M)=0\}
\]
はともに拡大に閉じる．
$\Phi(\mathcal Z_+)$はanticanonical heartの次数$0$に，
$\Phi(\mathcal Z_-)$はその$[-1]$ shiftに含まれる．
$V^0:\mathcal Z_+\to\operatorname{mod}_{fd}A$，
$V^1:\mathcal Z_-\to\operatorname{mod}_{fd}A$
はそれぞれ短完全列を保つ．
特に正指数（負指数）の既約根sheafの任意の有限回の
拡大は$\mathcal Z_+$（$\mathcal Z_-$）に入る．
\end{corollary}
\begin{proof}
$\Ext^1$，$\Hom$の各々の消滅は
\eqref{eq:offboundary-six}で拡大に関して閉じる．
定理\ref{thm:sign}により各既約根sheafは指定部分圏に属する．
長完全列の端の項が零なら残る三項は短完全列となる．
\end{proof}

\begin{proposition}[正負の混合拡大と相殺]
\label{prop:mixed-sign-cancellation}
$m_+>0$，$m_-<0$を持つ境界外既約根sheaf
$M_+\in\mathcal Z_+$，$M_-\in\mathcal Z_-$が同じ
Rains heartに実在すると仮定する．
$0\to M_-\to M\to M_+\to0$という拡大を取る．
この拡大が誘導する$A$-線型写像
\[
 \delta_M:V^0(M_+)\longrightarrow V^1(M_-)
\]
のrankを$r$とすると，
\begin{equation}\label{eq:mixed-ranks}
 \dim_kV^0(M)=m_+-r,\qquad
 \dim_kV^1(M)=-m_--r .
\end{equation}
特に$m_+=-m_->0$の場合，$\delta_M$が同型なら
$\Phi(M)=0$である．
従って$\Phi$が$\mathcal Z$上で保守的なら，
非零の混合拡大に対するそのような同型
$\delta_M$は存在しない．
同じ結論は$\Phi$が全導来圏の同値であるなら必ず成立する．
\end{proposition}
\begin{proof}
\eqref{eq:offboundary-six}に
$V^0(M_-)=V^1(M_+)=0$を代入すると
\[
0\to V^0(M)\to V^0(M_+)
\xrightarrow{\delta_M}V^1(M_-)\to V^1(M)\to0
\]
となり，\eqref{eq:mixed-ranks}が従う．
$\delta_M$が同型なら$N^0(M),N^1(M)$は
次数$0$の局所化が零であるが，$g$は各次数において
同型なので両方零であり，$\Phi(M)=0$である．
$M$は$M_-$を非零部分対象とするため非零である．
\end{proof}

\begin{remark}[Euler標数だけでは検出できない]
一般の$M\in\mathcal Z$に対して
\[
 \chi(e,M)=\dim V^0(M)-\dim V^1(M)
\]
であるが，$\chi(e,M)=0$から$\Phi(M)=0$は従わない．
正負の拡大類によるconnecting morphism
$\delta_M$のrankを指定して初めて
反標準切断のcohomologyが決まる．
$\mathcal Z_+,\mathcal Z_-$が$\mathcal Z$全体の
torsion pairになるとは証明していない．
\end{remark}

\begin{proposition}[根系からの共鳴交叉制約]
\label{prop:root-intersection}
Rains heartに境界に制限して零となる
異なる二つの既約純一次元根sheaf$M_\alpha,M_\beta$が
存在し，第一Chern類を$\alpha\ne\beta$とする．
すると
\[
 \Hom(M_\alpha,M_\beta)=\Ext^2(M_\alpha,M_\beta)=0,
 \qquad
 \dim_k\Ext^1(M_\alpha,M_\beta)=\alpha\cdot\beta.
\]
従って$\alpha\cdot\beta<0$ならそのような二対象は
同じRains heartには同時に存在できず，
$\alpha\cdot\beta=0$なら双方の向きの$\Ext^1$が零である．
$\alpha\cdot\beta=1$なら$\Ext^1$は一次元である．
\end{proposition}
\begin{proof}
非零な$M_\alpha\to M_\beta$の像$I$は，
$M_\beta$の純一次元性から一次元である．
$M_\alpha$の既約性より，核が一次元なら
その第一Chern類は$\alpha$となり，
像は零次元となって矛盾する．
核が零次元なら$M_\alpha$の純一次元性から零である．
従って写像は単射で，像の第一Chern類は$\alpha$である．
他方，$M_\beta$の既約性から非零部分sheafの第一Chern類は
$\beta$である．ゆえに$\alpha=\beta$となり矛盾する．

$\theta M_\alpha\simeq M_\alpha$であるから，
Serre双対により
$\Ext^2(M_\alpha,M_\beta)
 \simeq\Hom(M_\beta,\theta M_\alpha)^*=0$．
Rains\cite[Corollary 7.13, (7.4)]{RainsBirational}の
rank-zero Euler形式は
$\chi(M_\alpha,M_\beta)=-\alpha\cdot\beta$であり，
残る$\Ext^1$の次元は$\alpha\cdot\beta$となる．
\end{proof}

\begin{corollary}[導来完全忠実性に対する有限次元表現テスト]
\label{cor:affine-ext-test}
定理\ref{thm:offboundary-amplitude}の自然な導来切断$\Phi$が
二つの境界外既約根sheaf$M_+,M_-$（指数$m_+>0$，
$m_-<0$）の間の$\Ext^1$で完全忠実ならば，
対応する有限次元$A$加群$V_+$，$V_-$について
\[
\Ext^1_X(M_+,M_-)
 \xrightarrow{\;\sim\;}\Hom_A(V_+,V_-)
\]
が成り立つ．$\alpha_+=c_1(M_+)\ne c_1(M_-)=\alpha_-$なら
\[
\boxed{\dim_k\Hom_A(V_+,V_-)=\alpha_+\cdot\alpha_- .}
\]
従って右辺の次元がこれと異なるなら，
$\Phi$はその二対象の間で完全忠実ではない．
\end{corollary}
\begin{proof}
$\Phi(M_+)\simeq V_+$，$\Phi(M_-)\simeq V_-[-1]$
である（$V_\pm$は境界外の有限支持対象として見る）．
したがって
\[
\Hom_{\Db(\qgr R)}(\Phi M_+,\Phi M_-[1])
 =\Hom_{\qgr R}(V_+,V_-)=\Hom_A(V_+,V_-).
\]
最後の同型は，$g$が各$V_\pm$上で可逆であり，
その次数付き表示が$k[g]$に関して有限自由であることから，
局所化がこの二対象間の$\Hom$を保存するためである．
$\Ext^1_X(M_+,M_-)$は左辺へ$\Phi$が誘導する空間であり，
完全忠実性を仮定すると同型となる．
命題\ref{prop:root-intersection}を使えば次元式を得る．
この写像は命題\ref{prop:mixed-sign-cancellation}の
$\delta_M$のextension classによる線型写像そのものである．
\end{proof}

\begin{proposition}[根の和による反対指数の同時共鳴排除]
\label{prop:root-sum-exclusion}
古典的avatar$S$がgenuine del Pezzoで，
滑らかな反標準楕円曲線$E$を持つとする．
二つの根$\alpha,\beta$に
\[
 \mathcal O_E(\alpha)\otimes q^{m_\alpha}\simeq\mathcal O_E,
 \qquad
 \mathcal O_E(\beta)\otimes q^{m_\beta}\simeq\mathcal O_E
\]
という共鳴があるとする．もし$\alpha\cdot\beta=1$なら，
$m_\alpha+m_\beta=0$は不可能である．
特に，隣接する$A_2$部分根系の根に対する
指数$+r$と$-r$の共鳴は同時には起こらない．
\end{proposition}
\begin{proof}
$\gamma=\alpha+\beta$は
$\gamma^2=-2-2+2=-2$，$H\cdot\gamma=0$を
満たすから再び根である．
$m_\alpha+m_\beta=0$なら二つの共鳴式の積から
$\mathcal O_E(\gamma)\simeq\mathcal O_E$となる．
しかしgenuine del Pezzoについての
命題\ref{prop:zero-index}は任意の根の
$E$への制限が非自明であることを示している．
従って矛盾する．
\end{proof}

\begin{remark}[完全な相殺を探すときの制限]
異なる根sheafの間の非自明な拡大は，
上の命題から$\alpha\cdot\beta>0$を必要とする．
特に最も単純な$\alpha\cdot\beta=1$の場合，
指数が互いに反対で次元も等しいような混合拡大は，
genuine classical avatarではそもそも同時共鳴できない．
これは$\delta_M$の同型による導来的消失が
genuineな族に現れないことの\emph{完全な証明ではない}．
別のroot対や非既約対象を排除できていないためである．
\end{remark}

\begin{remark}[古典的な$(-2)$曲線上の完全な相殺]
\label{ex:euler-cancellation}
滑らかな弱del Pezzo $S$に$(-2)$曲線
$i:D\simeq\mathbb P^1\hookrightarrow S$があり，
反標準曲線$E$は$D$と交わらないとする．
$H=-K_S$の$D$への制限は自明である．
$\mathbb P^1$のEuler完全列
\[
0\to\mathcal O_D(-2)\xrightarrow{(-y,x)}
 \mathcal O_D(-1)^{\oplus2}
 \xrightarrow{(x,y)}\mathcal O_D\to0
\]
を$i_*$で送る．左端を$M_-$，右端を$M_+$
とすると$\chi(M_-)=-1$，$\chi(M_+)=1$である．
中間項$M=i_*\mathcal O_D(-1)^{\oplus2}$は非零だが
\[
\RHom_S(\mathcal O_S(nH),M)
 \simeq\mathbf R\Gamma(D,\mathcal O_D(-1)^{\oplus2})=0
 \quad(n\in\mathbb Z).
\]
したがって$\delta_M:k\to k$は同型で，
反標準導来切断は$M$を零に送る．
これは前の根壁の導来情報消失を
「正負指数の拡大の完全な相殺」として説明する
古典的な模型である．\emph{この例はweak avatarの例であり，
genuine非可換曲面で同じ拡大が存在するという主張ではない．}
\end{remark}


\section{正負の共鳴がそれぞれ稠密}
\begin{theorem}[正負各指数のZariski稠密性]
\label{thm:dense-both}
滑らかな平面三次曲線$E\subset\mathbb P^2$，
無限位数の$q\in\Pic^0(E)$を固定する．
$p_1,p_2,p_4,p_5,p_6\in E$を十分一般に固定し，
$L=\mathcal O_E(1)$と置く．各$n\ge1$と$\varepsilon\in\{+1,-1\}$に対し
\[
 \mathcal O_E(p_3^\varepsilon(n))\simeq
 L\otimes\mathcal O_E(-p_1-p_2)\otimes q^{-\varepsilon n}
\]
を満たす唯一の$p_3^\varepsilon(n)\in E$を取る．
各符号について有限個の例外を除いて六点は一般位置であり，
そのblow-upはgenuineで滑らかなcubic surfaceである．
単純根$\alpha=h-e_1-e_2-e_3$に対し
\[
 \mathcal O_E(\alpha)=q^{\varepsilon n},\qquad
 m=-\varepsilon n.
\]
Rainsの単純根sheaf存在定理の範囲では，
$\varepsilon=+1$の負指数sheafは
反標準切断の$[-1]$に移り，$\varepsilon=-1$の正指数sheafは
次数0にある．それぞれの符号の共鳴点集合が
Zariski denseである．
\end{theorem}
\begin{proof}
$\Pic^1(E)\simeq E$だから$p_3^\varepsilon(n)$は一意に存在する．
$q$が無限位数ゆえ各符号について$p_3^\varepsilon(n)$は相異なる．
五つの固定点に対し，追加の一点が三点共線性や六点
conic上条件を満たすのは有限個の点だけである．
従って一般位置を失う$n$は各符号について有限個であり，
$\{p_3^\varepsilon(n):n\ge1\}$はいずれも
$E$のZariski稠密部分集合である．
根の制限は$L(-p_1-p_2-p_3^\varepsilon(n))=q^{\varepsilon n}$．
式\eqref{eq:resonance}に$m=-\varepsilon n$を代入し，
定理\ref{thm:sign}を適用すれば結論を得る．
\end{proof}

\begin{corollary}[正負共鳴が全パラメータ空間で稠密]
\label{cor:dense-total}
滑らかな平面三次$E$，非torsionの$q\in\Pic^0(E)$，
および他の五吹き上げ点をパラメータとして動かす
六点配置の局所的な族において，定理\ref{thm:dense-both}の正指数および負指数の
共鳴点の合併はそれぞれZariski稠密である．
従ってこのチャートのいかなる非空Zariski開部分にも，
$\mathcal A_{\mathrm R}$で$(e,\tau)$がAZ-ampleでない負指数点が存在する．
\end{corollary}
\begin{proof}
各基底パラメータに対する$p_3$のファイバーは楕円曲線$E$であり，
無限位数$q$を持つ基底点はZariski稠密である．
それらの各ファイバーで定理の正負各指数列はZariski稠密なので，
全空間の任意の非空開集合は適当なファイバーでこの列に交わる．
基底全体のirreducibilityと$p_3$を動かす開写像性を用いる．
\end{proof}
\begin{remark}
ここでいう稠密性は可算個の閉共鳴軌跡の\emph{合併}の
Zariski稠密性であり，個々の共鳴軌跡が開集合であるとか，
全共鳴点で同じheart変化が起こるという主張ではない．
\end{remark}

\begin{remark}[他の次数への同じ稠密性]
$\mathbb P^2$の$r=3,\ldots,8$点blow-up（del Pezzo次数$d=9-r$）
で，滑らかな平面三次$E$上の他の点を一般に固定し，
$p_3(n)$を同じ式で動かせばよい．
genuine del Pezzo条件を破る$p_3$の集合は，固定した他の点が
十分一般なら$E$の有限集合である．
根$\alpha=h-e_1-e_2-e_3$はどの次数でも
$\alpha^2=-2$，$\alpha\cdot K=0$を満たすため，
Rainsの単純根sheaf存在結果の適用範囲で正負各指数の稠密性は
次数$1,\ldots,6$へ同様に及ぶ．これは全moduli成分への
稠密性を主張するものではない．
\end{remark}

\section{実在する混合$A_2$と二段拡大の非消失}\label{sec:dp-hearts2-mixed-a2}
本節の$E_6$最大共鳴の整数表示とZariski稠密性は主論文\cite{Hearts}を入力とする．
数値的共鳴だけから任意の根sheafの存在を推測せず，指定した単純根の実在を確認する．

\begin{proposition}[実在する正負の単純根sheaf]
\label{prop:actual-mixed-a2}
\[
 (N_1,\ldots,N_6)=(1,3,2,5,7,11),\qquad p_i=N_iq
\]
とする．六単純根のEuler指数は
\[
 (m_0,m_1,m_2,m_3,m_4,m_5)=(6,2,-1,3,2,4).
\]
特に
\[
 \alpha=e_1-e_2,\qquad\beta=e_2-e_3,\qquad
 m(\alpha)=2,\quad m(\beta)=-1,\quad
 \alpha\cdot\beta=1.
\]
Rains\cite[\S4, (4.14)--(4.15)]{RainsFiltered}の
simple-root effectivity判定と
\cite[\S9,\S11]{RainsBirational}の既約sheafの存在を
このsmooth elliptic chartに適用すると，同じRains heartに
境界非交叉の既約root sheaf
$M_\alpha,M_\beta$が存在する．それらについて
\[
 \dim\Ext^1(M_\alpha,M_\beta)=1,\quad
 \dim V^0(M_\alpha)=2,\quad
 \dim V^1(M_\beta)=1.
\]
またこの性質を保ち$N_1<N_3<N_2$とする全根共鳴点の集合も
同じ六点配置chartでZariski稠密である．
\end{proposition}
\begin{proof}
単純根指数は$h|_E=3o$と点の差から計算される．
全ての単純根の制限は$\langle q\rangle$にあるので，
Rainsのeffectivity判定によりこれらは有効である．
$E$は滑らかで$\alpha\cdot E=0$なので，
有効な単純根はformal $(-2)$-curveとして
境界と交わらない既約純一次元sheafで実現される．
無限位数の$q$のため$M_\alpha,M_\beta$のEuler指数は
上記の値に一意に固定される．
前節の命題\ref{prop:root-intersection}および
定理\ref{thm:sign}からExtと$V^j$の次元が従う．
最後は主論文\cite{Hearts}の全次数の最大共鳴と整数順序の稠密性による．
\end{proof}

\begin{proposition}[二段混合拡大全体での非消失]
\label{prop:a2-directsum}
命題\ref{prop:actual-mixed-a2}の二sheafを固定する．
$a\ge1,b\ge0$とし，任意の拡大
\[
0\longrightarrow M_\beta^{\oplus b}\longrightarrow M
 \longrightarrow M_\alpha^{\oplus a}\longrightarrow0
\]
を取る．$\Ext^1(M_\alpha,M_\beta)\simeq k$の生成元が
定める$A=R/(g-1)$加群写像を
$d:V_\alpha\to V_\beta$とする．
拡大類に対応する行列$T\in\operatorname{Mat}_{b\times a}(k)$の
connecting morphismは$d\otimes T$であり，
$r=\operatorname{rank}(d)\operatorname{rank}(T)$と置けば
\[
\boxed{
\dim V^0(M)=2a-r,\quad
\dim V^1(M)=b-r,\quad
r\le\min(a,b).}
\]
特に$a\ge1$なら$\dim V^0(M)\ge a>0$なので
自然な反標準導来切断はこの二段拡大を消さない．
$b=2a$では$\chi(e,M)=0$だが，依然として
$\Phi(M)\ne0$である．
\end{proposition}
\begin{proof}
直和とYoneda積の自然性により
\[
\Ext^1(M_\alpha^{\oplus a},M_\beta^{\oplus b})
 \cong\operatorname{Mat}_{b\times a}(k),
\]
対応するconnecting morphismは$k$-線型空間上
Kronecker積$d\otimes T$である．そのrankは積となり，
$d$は二次元から一次元への射だから$\operatorname{rank}(d)\le1$．
前節の六項完全列
\eqref{eq:offboundary-six}に代入して次元公式を得る．
\end{proof}

\begin{corollary}[全$q$導来同値に対する明示的な検査条件]
\label{cor:a2-ff}
同じ例で自然な導来切断$\Phi$が
$M_\alpha,M_\beta$の間の$\Ext^1$で完全忠実なら
\[
 \dim\Hom_A(V_\alpha,V_\beta)=1
\]
である．従って二次元表現$V_\alpha$は
一次元表現$V_\beta$を商に持ち，単純$A$加群ではない．
そのとき非分裂拡大のconnecting morphismはrank$1$で，
$(\dim V^0(M),\dim V^1(M))=(1,0)$となる．
\end{corollary}
\begin{proof}
前節の系\ref{cor:affine-ext-test}による．
$\dim V_\beta=1$だから非零$A$準同型は全射である．
$\Ext^1$の生成元は完全忠実性の下で非零の
$\Hom_A$へ写るから，connecting morphismのrankは$1$．
\end{proof}

\begin{remark}[有理数上の検算例]
\label{rem:a2-rational-witness}
\[
 E:\ y^2z=x^3-xz^2+z^3,\quad o=(0:1:0),\quad
 q=(0:1:1)
\]
とする．楕円曲線の二倍公式で
$2q=(\frac14:-\frac78:1)$である．
Nagell--Lutzの定理により$q$は無限位数である：
有限位数なら$2q$も有理torsionで整数座標を持つはずだが
$\frac14$は整数でない．
$N=(1,3,2,5,7,11)$の六点について
20通りの三点行列式と六点のconic評価行列式が
全て非零であり，72根の共鳴指数も全て非零であることを
\texttt{computations/verify\_resonance\_maximal\_a2.py}
により\emph{有理数の厳密演算}で検算した．
このコードは$A$の表現行列や$d$のrankを
計算するものではなく，Rainsのsheaf存在を
数値実験で置き換えてはいない．
\end{remark}

\begin{proposition}[全次数の最大共鳴]\label{thm:all-degree-maximal}
次数$d=1,\ldots,6$に対し，$p_i=N_iq$（相異なる正整数）なら，全根は非零指数で数値共鳴する．この事実と任意の整数順序を指定したZariski稠密性の完全な証明は主論文\cite{Hearts}にある．
\end{proposition}
\section{最大共鳴での既約根sheaf・zigzag Ext代数・二次障害}
\label{sec:resonance-zigzag}
本節では次数$d=1,\ldots,6$の最大共鳴点について，
「数値的に共鳴するすべての根が既約sheafを与える」
という誤った推論を排除し，実在するsheafのExt代数を
明示的に決定する．
先行研究からの入力は，Rains
\cite[\S4, pp.~33--34]{RainsFiltered}のsimple-root
effectivity判定と
\cite[Definition 9.1, Proposition 7.19, \S11]{RainsBirational}
の既約根sheafに関する結果である．
根の零でないEuler指数の整数表示は
定理\ref{thm:all-degree-maximal}により与えられている．

次数$d=9-r$，$3\le r\le8$とし，
単純根を
\[
\alpha_0=h-e_1-e_2-e_3,\qquad
\alpha_i=e_i-e_{i+1}\quad(1\le i<r)
\]
とする．単純根が実効的で，境界に制限して零の
既約純一次元sheaf $M_i$として実現される
Rains型smooth elliptic heartを固定する．
定理\ref{thm:all-degree-maximal}の
\emph{全根共鳴}の例ではこの仮定が
Rainsのsimple-root存在結果によって保証される．
$\Delta$を単純根のなすDynkin図とする．

\begin{theorem}[実在する既約根はちょうど単純根]
\label{thm:only-simple-irreducible}
$Li^*M=0$を満たす既約純一次元根sheaf $M$が
同じRains heartに存在するとき，
\[
 c_1(M)\in\{\alpha_0,\ldots,\alpha_{r-1}\}
\]
である．逆に各$\alpha_i$は上の$M_i$で実現される．
$q$が非torsionならば，境界非交叉の既約根sheafは
各$\alpha_i$について同型を除き\emph{一つだけ}である．
従って，数値的に共鳴している根の総数と
実際の既約根sheafの同型類数は
\[
\begin{array}{c|rrrrrr}
d&6&5&4&3&2&1\\\hline
\#\text{数値的共鳴根}&8&20&40&72&126&240\\
\#\text{既約根sheaf}&3&4&5&6&7&8
\end{array}
\]
と異なる．
\end{theorem}
\begin{proof}
根系の正根$\gamma$が単純根でないとする．
$\gamma=\sum_{i=0}^{r-1}c_i\alpha_i$，
$c_i\in\mathbb Z_{\ge0}$と書くと
\[
 -2=\gamma^2=\sum_i c_i(\gamma\cdot\alpha_i)
\]
なので，ある$c_i>0$について$\gamma\cdot\alpha_i<0$である．
しかし既存の命題\ref{prop:root-intersection}によれば，
異なる根類の境界非交叉既約根sheafが同時に存在するなら
$\dim\Ext^1(M_\gamma,M_i)=\gamma\cdot\alpha_i\ge0$
でなければならない．従って$\gamma$は不可能．

負根$\gamma$については$-\gamma$が正根であり，
有効な$\alpha_i$の非負整数和なので$-\gamma$は有効である．
もし$M_\gamma$が存在すれば$\gamma$も有効だから，
Rains\cite[Proposition 7.19]{RainsBirational}の
「$D$と$-D$がともに有効なら$D=0$」に反する．
従って単純正根だけが残る．

一意性を証明する．$c_1(M)=c_1(N)=\alpha_i$で
$Li^*M=Li^*N=0$なら，
restriction determinantの式\eqref{eq:resonance}と
$q$の非torsion性から
$\chi(e,M)=\chi(e,N)$である．
従って$[M]=[N]$ in $K_0^{\mathrm{num}}$で，
Rains\cite[Corollary 7.13]{RainsBirational}により
$\chi(M,N)=-\alpha_i^2=2$．
$\theta M\simeq M,\theta N\simeq N$なので，
もし両向きの$\Hom$が零ならSerre双対性で両向きの
$\Ext^2$も零となり，$\chi(M,N)=-\dim\Ext^1(M,N)\le0$
となって矛盾する．従っていずれかの向きに非零射がある．
既約性と純一次元性により，非零射のkernelは零である：
もし非零kernelがあれば同じChern類を持ち，
像は零次元となってtargetの純一次元性に反する．
両者の数値類が同じなので，
Rains\cite[Corollary 7.7]{RainsBirational}から
単射は同型である．
\end{proof}

\begin{theorem}[Dynkin zigzag Ext代数]
\label{thm:zigzag-ext}
$T=\bigoplus_{i=0}^{r-1}M_i$と置く．
$\operatorname{Ext}^{\bullet}_{\mathcal A_{\mathrm R}}(T,T)$は
次の明示的なgraded algebra $Z_\Delta$と同型である．
\begin{enumerate}
\item 次数$0$には各頂点の冪等元$1_i$が一個ずつある．
\item 次数$1$には$\Delta$の各辺$i-j$に対して
両向きの矢$a_{ij}$，$a_{ji}$が各一個ある．
次数$2$には各頂点に一個のloop $z_i$がある．
\item 異なる頂点へ向かう長さ二の道の積は零．
各辺$i-j$の往復積は
$a_{ij}a_{ji}=c_{ij}z_i$，
$a_{ji}a_{ij}=c_{ji}z_j$（$c_{ij},c_{ji}\ne0$）である．
頂点基底と矢をrescaleすれば通常のzigzag関係式にできる．
\item 次数$>2$の積はすべて零．孤立頂点（$A_1$成分）
も次数$2$の$z_i$を持つ．
\end{enumerate}
従って$\#V(\Delta)=r$，$\#E(\Delta)=s$とすると
\[
\dim_k\operatorname{Ext}^{\bullet}(T,T)=2r+2s
=\begin{cases}
8,&d=6,\\
14,18,22,26,30,&d=5,4,3,2,1\text{ の順}.
\end{cases}
\]
これはまず\emph{次数付きYoneda積}の同定である．
dgおよびminimal $A_\infty$模型の形式性については，
以下でLiu--Wangの既知定理を用いて判定する．
\end{theorem}
\begin{proof}
定理\ref{thm:sign}と命題\ref{prop:spherical}によれば
$\Ext^0(M_i,M_i)=\Ext^2(M_i,M_i)=k$，
$\Ext^1(M_i,M_i)=0$．
異なる$i,j$については
命題\ref{prop:root-intersection}より
\[
\Ext^0(M_i,M_j)=\Ext^2(M_i,M_j)=0,\qquad
\dim\Ext^1(M_i,M_j)=\alpha_i\cdot\alpha_j,
\]
これは隣接頂点で$1$，それ以外で$0$である．
高次Extは滑らかなRains heartのhomological dimension $2$
から消える．$\theta\to\mathrm{id}$は全$M_i$上で同型だから，
これらが生成する部分三角圏ではSerre関手が自然に$[2]$となる．
よってSerre traceは
\[
\Ext^1(M_i,M_j)\otimes\Ext^1(M_j,M_i)
\longrightarrow \Ext^2(M_i,M_i)\xrightarrow{\operatorname{tr}} k
\]
を完全な双線型pairingにする．
各辺で往復積が非零となり，
二つの異なる頂点間の$\Ext^2$消滅から
非往復の長さ二の積は零である．
全頂点に次数二の一次元空間が一つずつあるため，
$\Delta$が木（$A_2+A_1$では森）であることを用い
矢とloopの基底をrescaleすれば
各頂点で往復積が同じloopとなる．
以上がzigzag代数の表示である．次数別の次元は
$r,2s,r$なので合計$2r+2s$となる．
\end{proof}

\begin{corollary}[二次Kuranishi障害のpreprojective型表示]
\label{cor:zigzag-quadratic}
$T=\bigoplus_iM_i$の一次変形を
$x=\sum_{i-j}x_{ij}a_{ij}\in\Ext^1(T,T)$
と書く．その二次Kuranishi障害
$m_2(x,x)\in\Ext^2(T,T)$は，適切な非零定数$c_{ij}$を用い
\[
m_2(x,x)=\sum_i
\left(\sum_{j\sim i}c_{ij}x_{ij}x_{ji}\right)z_i
\]
である．$\Delta$が森なので，この二次障害の消滅条件は
\[
\boxed{x_{ij}x_{ji}=0\qquad(\{i,j\}\in E(\Delta))}
\]
と\emph{イデアルとしても}同値である．
従って二次障害方程式で定まるスキーム（gauge商を取る前）は
$2^s$個の座標線型空間の合併で，その各々は次元$s$である．
$d=3$の$E_6$なら$s=5$なので$32$成分となる．
ここまでは一次Yoneda積から得る二次方程式の計算である．
高次障害が追加されないことは次の定理によって従う．
\end{corollary}
\begin{proof}
次数二のYoneda積は前定理の往復積以外は零なので
最初の表示を得る．
任意の葉$i$について次数二の係数方程式は
$c_{ij}x_{ij}x_{ji}=0$ただ一つであり$c_{ij}\ne0$．
従ってその辺積は障害イデアルに属する．
この辺の項を隣の頂点$j$の方程式から取り去って
森の残りのグラフで帰納的に繰り返す．
よって障害イデアルは
$(x_{ij}x_{ji}:\{i,j\}\in E(\Delta))$に等しい．
これは互いに異なる変数対を使う平方自由単項式のイデアルで，
各辺から一方の矢を零にする$2^s$個の座標部分空間を
既約成分とし，各成分では残りの矢が$s$個自由である．
\end{proof}

\begin{theorem}[根配置のdg形式性と形式変形の二次完全性]
\label{thm:zigzag-formal}
基礎体が標数$0$で，$\Delta$が上の
$A_2+A_1,A_4,D_5,E_6,E_7,E_8$のいずれかなら，
$T=\bigoplus_i M_i$の導来自己準同型dg代数は
\[
\boxed{\mathbf R\operatorname{End}_{\mathcal A_{\mathrm R}}(T)
 \ \simeq_{A_\infty}\ Z_\Delta}
\]
として形式的である．従って\emph{固定した}Rainsの
enhanced導来圏の中で$T$の形式的変形を考え，
各直和因子の標識を固定したMaurer--Cartanの
framedパラメータ空間では，
\[
\boxed{
 \widehat{\mathcal O}_{\rm MC}
 \cong
 k[\![x_{ij},x_{ji}:\{i,j\}\in E(\Delta)]\!]/
      (x_{ij}x_{ji}:\{i,j\}\in E(\Delta)).
}
\]
よって高次$A_\infty$障害による追加の切断はなく，
この形式的スキームは$2^{s}$個の
$s$次元formal座標部分空間の合併である
（$s=\#E(\Delta)$）．
特にcubic $E_6$では$32$個の$5$次元formal分枝を持つ．
\end{theorem}
\begin{proof}
Liu--Wang\cite[Theorem~1.1]{LiuWangZigzag}は，
標数$0$における連結ADE Dynkin木の
\emph{path-length graded} zigzag代数が
intrinsically formalであることを示す．
ここで$Z_\Delta$のpath lengthは先のExtの
cohomological degree $0,1,2$に一致している．
Kadeishviliのminimal模型定理により
$\mathbf R\operatorname{End}(T)$の最小$A_\infty$模型は
その$\Ext^\bullet$上に取れるので，
前定理とLiu--Wangを合わせると全ての高次
$m_n$ $(n\ge3)$は$A_\infty$準同型で消せる．

$d=6$の$A_2+A_1$では両連結成分間の
全$\Ext^\bullet$が零なので，minimal模型は
成分ごとに直和へ分離する．
$A_2$にはLiu--Wangを使い，孤立$A_1$の
$\Ext^\bullet=k\oplus kz$（$\deg z=2$，$z^2=0$）
については，strictly unitalな$m_n$ $(n\ge3)$は
正の次数の入力しか持たないと出力次数が
$2n+2-n=n+2>2$になり，次数の理由で零となる．
よってこの場合も形式的である．

Framedな変形のMaurer--Cartan元は
$\Ext^1(T,T)\widehat\otimes\mathfrak m$に属する
$x=\sum x_{ij}a_{ij}$であり，
形式性から方程式は$m_2(x,x)=0$だけで足りる．
前系の葉からの帰納法でその定義イデアルは
$(x_{ij}x_{ji})$そのものとなる．
従って表示した完成環とformal分枝数を得る．
各直和因子の標識を忘れた変形の群oidは
対角の冪等元に対応するgauge群の作用で
さらに商を取るべきであり，ここでは扱わない．
\end{proof}
\begin{remark}[形式性が与えないもの]
この定理はRainsの\emph{固定された}dg圏における
$T$の形式的変形の記述である．
Rains曲面や楕円境界のパラメータを動かす
相対変形，anticanonical heartの全体，
$A=R/(g-1)$の有限次元表現のモジュライ，
および混合根のconnecting map$\delta$の
具体的rankは，この形式性からは決定しない．
\end{remark}

\begin{remark}[導来比較を監査する新しい有限モデル]
特にcubicの最大共鳴で実在する$6$個の単純根sheafの
graded Yoneda algebraは次元$22$の$E_6$ zigzagである．
Rains--anticanonicalの指定した導来切断が完全忠実かを調べるには，
このgraded algebraの積だけでなく，
$\mathbf R\operatorname{Hom}(e,T)$の$A_\infty$-module構造，
とりわけ正負隣接根間の
$\operatorname{Ext}^1$から誘導される
$V_\alpha\to V_\beta$のconnecting morphismを決定する必要がある．
ここでの形式性は上の定理で確立したが，
このことだけでは$\delta$のrankは決まらない．
具体的な$A$加群の作用行列とnatural functorの
$A_\infty$-module構造の計算が独立に必要である．
\end{remark}

\section{非torsion境界外圏と混合\texorpdfstring{$A_2$}{A2}の完全判定}
\label{sec:mixed-a2-category}
本節は自然な導来切断の完全忠実性を仮定せず，混合根のconnecting mapを決定する．
Rains\cite[Corollaries 7.7, 7.13, 7.15, Propositions 7.6, 7.32,
Corollary 10.12, Proposition 10.13]{RainsBirational}と既存の符号別Extを入力とする．
$H=-K_X$はnef，$H^2=d>0$，$q$は非torsion，構造line bundle $e$はexceptionalとする．
$\Ext_X$は常にambientのRains heartで計算する．

\subsection{境界外圏の有限長性}
\begin{lemma}[零次元の境界外商の排除]
\label{lem:offboundary-no-points}
$Li^*M=0$なら，$M$の零次元商は零である．従って境界非交叉の
既約純一次元sheafはRains heartの単純対象である．
\end{lemma}
\begin{proof}
零次元商$M\twoheadrightarrow Z$には右完全性から$i^*Z=0$が従い，
$\theta Z\to Z$は全射となる．両者の数値類は等しいので
Rains Corollary 7.7により同型，従って$Li^*Z=0$である．
一方$[Z]=\ell[\mathrm{pt}]$，$\ell=\chi(e,Z)\ge0$で，
境界restrictionのdeterminantは$q^\ell$となる
\cite[Proposition 7.6, Proposition 7.24直前]{RainsBirational}．
非torsion性から$\ell=0$，ゆえに$Z=0$．既約純一次元sheafの
非零部分対象の商は零次元なので単純性が従う．
これは非torsion性と境界非交叉を使う証明であり，既約性だけから単純性を推論してはいない．
\end{proof}

\begin{proposition}[境界外Serre部分圏]
\label{prop:offboundary-serre-length}
$\mathcal Z=\{M:\dim M\le1,\ Li^*M=0\}$は有限長Serre部分圏であり，
その単純対象は2-sphericalな境界非交叉根sheafである．
最大共鳴で全単純根sheafが実在する場合，単純対象は
定理\ref{thm:only-simple-irreducible}の$r=9-d$個に尽きる．
\end{proposition}
\begin{proof}
$0\to N\to M\to Q\to0$，$M\in\mathcal Z$とする．
nef性と加法性から$0=Hc_1(M)=Hc_1(N)+Hc_1(Q)$の両項は零．
$i^*Q=0$なので$\theta Q\to Q$は全射であり，rank零と$K_Xc_1(Q)=0$から
両者の数値類は同じ\cite[Corollary 7.15]{RainsBirational}．
Corollary 7.7から同型となり，$Li^*Q=0$．自然性から$\theta N\to N$も同型．
拡大閉性は境界triangleによる．
非零零次元対象がないため，真の降鎖の逐次商はすべて一次元である．
Rains Proposition 7.32のfinitely partitive性とnoetherian性から有限長となる．
単純対象$S$の$D=c_1(S)\ne0$は$HD=0$を満たす．
$H^\perp$の負定値性，Riemann--Rochの偶奇，$\theta S=S$とSerre双対性より
\[
0<-D^2=\chi(S,S)=2-\dim\Ext_X^1(S,S)\le2
\]
であり，$D^2$は偶数だから$D^2=-2$，$\Ext_X^1(S,S)=0$．
最後は既存の単純根の排他性と一意性を使う．
\end{proof}

\subsection{任意の正負指数に対する最大rank}
隣接する境界非交叉既約根sheafを$S_+,S_-$とし，
\[
\alpha=c_1(S_+),\quad\beta=c_1(S_-),\quad\alpha\beta=1,
\qquad\chi(e,S_+)=p>0,\quad\chi(e,S_-)=-n<0
\]
とする．$\gamma=\alpha+\beta$は根である．非分裂拡大を
\begin{equation}\label{eq:a2-two-extensions}
0\to S_-\to P\to S_+\to0,\qquad
0\to S_+\to Q\to S_-\to0
\end{equation}
と書く．両向きのExt$^1$は一次元なので各中間項は同型を除き一意である．

\begin{theorem}[任意の正負指数のconnecting rank]
\label{thm:a2-delta-rank}
$P$のconnecting map$\delta:k^p\to k^n$は
\[
\boxed{\operatorname{rank}\delta=\min(p,n),\qquad
 (\dim V^0(P),\dim V^1(P))=((p-n)_+,(n-p)_+)}
\]
を満たす．ここで$x_+=\max(x,0)$．
$p>n$なら非零$e\to P$は全射で，核$L$はrank-one torsion-free，
\[
c_1(L)=-\gamma,\quad\chi(e,L)=1-p+n,\quad
\chi(L,L)=3-2(p-n).
\]
$L$がline bundleとなることと$p-n=1$は同値である．
特に$p=2,n=1$ではrankは1，$(h^0,h^1)=(1,0)$である．
いずれにも導来完全忠実性を仮定しない．
\end{theorem}
\begin{proof}
符号別Extから$\Hom(e,P)=\ker\delta$，$\Ext^1(e,P)=\operatorname{coker}\delta$．
非零$f:e\to P$があるとする．$\Hom(e,S_-)=0$なので$P\to S_+$との合成は非零．
$S_+$の単純性から全射である．像$I$と$S_-$の交わりは$0$か$S_-$であり，
前者なら拡大を分裂させてしまう．従って$I=P$，すなわち$f$は全射．
その核$L$はrank-one torsion-freeで$c_1(L)=-\gamma$である．
Serre双対性とRains Corollary 10.12より
\[
\Ext^2(e,L)=\Hom(L,\theta e)^*=0:
\quad H(c_1(\theta e)-c_1(L))=H(K_X+\gamma)=-d<0.
\]
ここでCorollary 10.12はrank-one torsion-free同士に適用でき，
$L$をline bundleと仮定する必要はない．
$e$のexceptionalityと$0\to L\to e\to P\to0$から$\Ext^1(e,P)=0$となる．

$p>n$なら$\ker\delta\ne0$なので上の議論が適用でき，$\delta$は全射．
$p\le n$なら$\Hom(e,P)\ne0$とすると$h^1(P)=0$になり，
$0\ge p-n=\chi(e,P)=h^0(P)>0$となって矛盾する．
従ってこの場合は$\delta$が単射．これで最大rankを得る．
$p>n$の場合，Euler加法性とRiemann--Rochから表示した$L$の値を得る．
Rains Proposition 10.13の等号条件$\chi(L,L)=1$により，
line bundle条件は$p-n=1$と同値である．
\end{proof}
\begin{remark}[指数の釣合いと消失]
$p=n$で上の二sheafが存在すれば，$P\ne0$なのに$\Phi(P)=0$となる．
しかしgenuine smooth classical avatarでは$\gamma$が零指数の共鳴根となるため，
この釣合いの場合自体が起こらない．一般の数値条件からこのsheafの実在を主張しない．
\end{remark}

\begin{corollary}[二段直和拡大]
\label{cor:a2-rank-matrix-exact}
$0\to S_-^{\oplus b}\to M\to S_+^{\oplus a}\to0$の拡大類を
$T\in\operatorname{Mat}_{b\times a}(k)$，$t=\operatorname{rank}T$，$r=\min(p,n)$とすると
\[
\operatorname{rank}\delta_M=rt,\qquad
\dim V^0(M)=pa-rt,\qquad\dim V^1(M)=nb-rt.
\]
$p=2,n=1$では$(2a-t,b-t)$となり，$a\ge1$ならEuler標数が零でも消失しない．
\end{corollary}
\begin{proof}
Yoneda積の自然性から$\delta_M=\delta\otimes T$であり，rankは積となる．
六項完全列を適用する．
\end{proof}

\subsection{有限長部分圏の全対象と局所torsion pair}
\begin{theorem}[四つの直既約対象]
\label{thm:a2-serre-category}
$\mathcal C=\langle S_+,S_-\rangle_{\rm Serre}$は
$\operatorname{mod}_{\rm fd}\Pi(A_2)$と同値であり，直既約は$S_+,S_-,P,Q$だけである．
$P\oplus Q$はprojective generatorで，$\Pi(A_2)$は二頂点・両向き二矢と
両方の長さ二の道を零とする関係を持つ次元4の代数である．
\end{theorem}
\begin{proof}
非分裂性と単純性から，順序$(S_+,S_-,P,Q)$で
\[
(\dim\Hom(U,V))_{U,V}=
\begin{pmatrix}1&0&0&1\\0&1&1&0\\1&0&1&1\\0&1&1&1\end{pmatrix}.
\]
例えば$\Hom(S_+,P)=\Hom(P,S_-)=0$は，connecting map
$k\to\Ext_X^1(S_+,S_-)=k$が非零拡大類を送ることから従う．
$\Hom(P,Q)=k$は$P\twoheadrightarrow S_+\hookrightarrow Q$で生成され，
逆方向も同様で，両合成は零である．
$\theta$が全対象で自然に恒等なので
\[
\dim\Ext_X^1(U,V)=\dim\Hom(U,V)+\dim\Hom(V,U)+c_1(U)c_1(V).
\]
上の表を代入すると，$S_+,S_-$間の両向きの一次元Ext以外は全て零となる．
とくに$P,Q$から両単純へのExt$^1$は零．長さによる帰納法から
$P,Q$は$\mathcal C$でprojectiveで，単純への全射を持つのでgeneratorとなる．
そのEnd代数のHomと合成が上記の表示を与える．
表現を$x:U_+\to U_-$，$y:U_-\to U_+$，$xy=yx=0$で書くと，
$\operatorname{im}y\subset\ker x$，$\operatorname{im}x\subset\ker y$に適合する
基底により，一方の矢のみ非零の一次元対と孤立一次元空間へ分解される．
これが四つの直既約である．
\end{proof}

\begin{corollary}[全cohomologyと指数に応じた局所torsion pair]
\label{cor:a2-torsion-pair}
$\mathcal C$の直既約の$(h^0,h^1)$は
\[
S_+:(p,0),\quad S_-:(0,n),\quad P:((p-n)_+,(n-p)_+),\quad Q:(p,n).
\]
一般のquiver表現で$a=\dim U_+$，$b=\dim U_-$，$t=\operatorname{rank}x$なら
\[
h^0=pa-\min(p,n)t,\qquad h^1=nb-\min(p,n)t.
\]
$p>n$のとき$(\ker h^1,\ker h^0)=(\operatorname{add}(S_+,P),\operatorname{add}(S_-))$，
$p<n$のとき$(\operatorname{add}(S_+),\operatorname{add}(S_-,P))$はtorsion pairである．
$p\ne n$なら$\Phi$は非零対象を消さない．$p=n$の場合の核は$\operatorname{add}(P)$で，
二つのkernel部分圏はtorsion pairにならない．
\end{corollary}
\begin{proof}
$P$は最大rank定理，$Q$は逆向きの六項完全列から求まる．
四つの直既約の重複度は順に
$a-\operatorname{rank}x-\operatorname{rank}y$，
$b-\operatorname{rank}x-\operatorname{rank}y$，
$\operatorname{rank}x$，$\operatorname{rank}y$であり，直和して公式を得る．
Hom表がtorsion pairの直交性を与え，$Q$の短完全列がtorsion分解を与える．
他の直既約は片方の部分圏に属する．非消失および核の記述もこの表から従う．
\end{proof}
\begin{remark}
$\chi(e,P)=\chi(e,Q)=p-n$だが，$Q$には二つのcohomological degreeが残る．
非既約対象ではEuler指数の符号だけからheart内の位置を決定できない．
また$\Ext_{\mathcal C}^{2j}(S_+,S_+)=k$ $(j\ge0)$とambientの
$\Ext_X^{>2}=0$は異なる．従って$D^b(\mathcal C)\to D^b(\mathcal A_q)$は
完全忠実でなく，このabel分類をambient三角圏の同定に読み替えない．
\end{remark}

\subsection{二次元表現の検査条件と導来比較への接続}
$p=2,n=1$に戻り，$V_+=V^0(S_+)$，$V_-=V^1(S_-)$，$W=V^0(P)$とする．
\begin{proposition}[単一Homテストの正確な範囲]
\label{prop:a2-scalar-test}
$0\to W\to V_+\xrightarrow{\delta}V_-\to0$は$A$加群の短完全列で，両端は一次元．
$\dim\Hom_A(V_+,V_-)$は$V_+\simeq V_-^{\oplus2}$の場合のみ2，それ以外は1である．
自然写像$\Ext_X^1(S_+,S_-)\to\Hom_A(V_+,V_-)$は常に単射である．
\end{proposition}
\begin{proof}
六項完全列と最大rankから短完全列を得る．準同型空間は
二次元の$\Hom_k(V_+,V_-)$の非零部分空間．独立な二準同型があれば
それらを並べて$V_+\simeq V_-^2$を得る．拡大類の生成元が非零$\delta$に写るので単射性も従う．
\end{proof}

右加群を行ベクトルで表し，$W$を第一座標に取ると
\[
\rho(a)=\begin{pmatrix}\chi_W(a)&0\\\eta(a)&\chi_-(a)\end{pmatrix},\qquad
\eta(ab)=\eta(a)\chi_W(b)+\chi_-(a)\eta(b).
\]
\begin{proposition}[導来完全忠実性の必要条件]
\label{prop:a2-character-obstructions}
$\Phi$が二根の生成する三角圏上で完全忠実なら，$W\not\simeq V_-$で，
上の加群拡大は非分裂である．従って作用像は全三角行列代数である．
\end{proposition}
\begin{proof}
$\Ext_X^1(P,S_-)=0$から完全忠実性の下で$\Hom_A(W,V_-)=0$．
また非零包含$S_-\hookrightarrow P$は加群拡大の非零類へ移る．
異なるcharacterと非分裂性により二つの対角冪等元と非零非対角成分を得る．
\end{proof}
\begin{remark}
同じcharacterのJordan型は単一Hom次元1のテストには合格しても上の必要条件に反する．
異なるcharacterでも分裂型なら非零sheaf射を消す．対象の非消失と射の忠実性は別である．
ただし本稿のgenuine楕円gluingでは，次節の比較定理と
定理\ref{thm:resonance-affine-triangular}により両方の必要条件が実際に成立する．
残るのは元の$A$の生成元に対するcharacterと非対角係数の具体的な式であって，
この範囲での分離・非分裂性の真偽ではない．
\end{remark}



\section{零指数とgenuine性の境界}\label{sec:resonance-zero-summary}
\begin{proposition}[genuine avatarにおける零指数の排除]\label{prop:zero-index}
滑らかなgenuine del Pezzo $S$と滑らかな反標準曲線$E$について，任意の根$\alpha$は$\mathcal O_E(\alpha)\ne\mathcal O_E$を満たす．従って非torsionな$q$に関して零指数根共鳴はない．
\end{proposition}
\begin{proof}
$-K_S$のamplenessと根のRiemann--Rochから$H^0(S,\pm\alpha)=0$であり，制限完全列から$H^0(E,\mathcal O_E(\alpha))=0$を得る．詳細は主論文\cite{Hearts}を参照する．
\end{proof}
零指数壁の古典的収縮・非可換弱gluingのorbit核とdg商は主論文で扱う．

\section{再現と今後の重点}
符号別Euler計算と根格子は\texttt{computations/verify\_resonance\_hearts.py}，
拡大の相殺と根の交叉は\texttt{computations/verify\_resonance\_extensions.py}，
全根共鳴の有理楕円点・全次数の根数・整数指数は
\texttt{computations/verify\_resonance\_maximal\_a2.py}，
既約根の排除条件とzigzag分枝数は\texttt{computations/verify\_resonance\_zigzag.py}で検算する．
混合$A_2$のHom/Euler/Ext，任意正負指数の両chamber，四次元商の行列模型，
および局所Weyl反射は\texttt{computations/verify\_resonance\_mixed\_a2.py}で検算する．
これらはsheaf存在や導来比較を計算で証明するものではない．

現在の具体的な未解決問題は，元の楕円データから$A=R/(g-1)$の生成元を固定し，
商$A\twoheadrightarrow\Pi(A_2)$の各係数を明示することである．
characterの分離，二次元拡大の非分裂性，および指定gluingでの全heartのtorsion pairは
もはや未決定として扱わない．一方，全$E_6$普遍有限商での楕円生成元の具体的作用行列，独立のavatar間での指定同定，相対family比較と退化境界は引き続き課題である．
実行済み検算・組版・文献照合の正確な範囲は
\texttt{research/resonance\_hearts\_validation.json}に記す．

